% Generated draft: requires mathematical and source review.
\documentclass[11pt]{article}
\usepackage[margin=25mm]{geometry}
\usepackage{fontspec}
\setmainfont{Latin Modern Roman}
\usepackage{amsmath,amssymb,mathtools,unicode-math}
\setmathfont{Latin Modern Math}
\usepackage{enumitem,needspace,xurl,hyperref}
\setlength{\parindent}{0pt}
\setlength{\parskip}{5pt}
\setlength{\emergencystretch}{3em}
\newcommand{\N}{\mathbb{N}}
\newcommand{\Z}{\mathbb{Z}}
\newcommand{\Q}{\mathbb{Q}}
\newcommand{\R}{\mathbb{R}}
\newcommand{\C}{\mathbb{C}}
\newcommand{\F}{\mathbb{F}}
\newcommand{\A}{\mathbb{A}}
\newcommand{\PP}{\mathbb{P}}
\newcommand{\E}{\mathbb{E}}
\newcommand{\eps}{\varepsilon}
\newcommand{\dd}{\,\mathrm{d}}
\begin{document}
\section*{UnsolvedMath ID: 20002481}
\textbf{Source problem code:} catalogue entry
\subsection{Definitions and mathematical statement}
This record merges two distinct AIM entries, Problems 23 and 24 in the supplied source-verification account. Both questions are retained.

\textbf{Problem 23: a priori stable Arnold diffusion.} Let \(\mathbb T=\R/\Z\), \(n\ge2\), and
\[
 H_\varepsilon(x,p,t)=H_0(p)+\varepsilon H_1(x,p,t),
 \qquad (x,p,t)\in\mathbb T^n\times\R^n\times\mathbb T,
 \quad 3\le r\le\omega.
\]
The Hamiltonian is \(C^r\); \(r=\omega\) denotes real analytic regularity. The source assumes \(H_0\) positive definite and superlinear; interpreting positive definite as positive definite Hessian gives \(D^2H_0(p)>0\), and superlinearity means \(H_0(p)/|p|\to+\infty\) as \(|p|\to\infty\). The precise regularity topology on this noncompact phase space is not specified.

For given nonempty open sets \(U_0,U_1\subset\R^n\), the conjectured property is that a typical perturbation \(H_1\), fixed independently of \(\varepsilon\), satisfies
\[
 \exists\varepsilon_0>0\ \forall\varepsilon\in(0,\varepsilon_0)\
 \exists (x_\varepsilon(t),p_\varepsilon(t)),\ t_0<t_1:
 \quad \dot x_\varepsilon=\partial_pH_\varepsilon,\quad
 \dot p_\varepsilon=-\partial_xH_\varepsilon,\quad
 p_\varepsilon(t_i)\in U_i\ (i=0,1).
\]
The orbit and visiting times may depend on \(\varepsilon\); no diffusion-time bound is asserted. The source explicitly leaves the meaning of typical to be established. It reports an \(n=2\), \(C^\infty\), cusp-residual theorem of Mather for perturbations \(\varepsilon H_1\), while stressing that its allowed amplitudes along a fixed direction need not contain a terminal interval \((0,\varepsilon_0)\). It asks whether an analogue holds in a cusp-residual class of high-order trigonometric polynomials of order \(N\), allowing exponentially small coefficient changes. The source cites exponentially slow analytic action change under Nekhoroshev theory as an obstacle and suggests finite regularity may be easier; the relevant stability hypotheses and estimates are not specified.

\textbf{Problem 24: a priori unstable methods.} On \(\mathbb T^n\times\R^n\times\R^2\), consider the autonomous Hamiltonian
\[
 H(\theta,I,x,y)=H_0(I)+F(x,y)+\varepsilon H_1(\theta,I,x,y),
\]
where \(H_0\) is convex and superlinear, \(F:\R^2\to\R\) has a separatrix loop, and \(H_1\) is generic in a class not specified. Understand the relation between the variational method and the separatrix-map method for overcoming the large-gap problem of Arnold diffusion. The source credits variational methods of Mather (Xia), geometric methods using secondary KAM tori to de la Llave, Delshams and Seara, and the separatrix-map method to Treschev. Their exact theorem hypotheses are not given. This comparison question does not assert arbitrary-endpoint diffusion across different energy levels of the autonomous system.

\textbf{Scope of the supplied auxiliary aid.} Ambient genericity, even with open dense sets of good amplitudes along rays, does not by itself imply the all-small-amplitudes quantifier above. Nor does an ambient residual or open dense class automatically give genericity in a prescribed Fourier family. A fixed Fourier cutoff in periodic variables can still leave infinitely many action-dependent coefficient functions; the aid's finite-dimensional argument additionally requires a chosen coefficient parameterization.

The aid gives the following conditional finite-dimensional reduction. Let \(a\in\Omega\subset\R^d\), with \(\Omega\) open, and suppose finitely many sufficient certification conditions for a specified diffusion construction are \(D_j(\varepsilon,a)\ne0\), where
\[
 D_j(\varepsilon,a)=\varepsilon^{m_j}\bigl(d_j(a)+o(1)\bigr),
 \quad\varepsilon\downarrow0,
\]
locally uniformly in \(a\). Assume each \(d_j\) is real analytic and not identically zero on any connected component of \(\Omega\). Then the aid reports that
\[
 G=\{a\in\Omega:\ d_j(a)\ne0\text{ for every }j\}
\]
is open dense, and every \(a\in G\) satisfies all certification conditions for every sufficiently small positive amplitude, with a common cutoff on each compact subset of \(G\). This does not establish such finite certifications for the original Hamiltonian problem; flat or exponentially small conditions need not have this leading-power form.

For the method comparison, a local generating function \(S(q,Q)\) represents an exact symplectic map by
\[
 p=-\partial_1S(q,Q),\qquad P=\partial_2S(q,Q).
\]
With endpoints fixed, critical sequences of \(\sum_k S(q_k,q_{k+1})\) satisfy
\[
 \partial_2S(q_{k-1},q_k)+\partial_1S(q_k,q_{k+1})=0,
\]
which is precisely momentum matching for consecutive map steps, on a domain where the generating function defines the map. This local identity is the aid's bridge between variational sequences and twist-kick separatrix maps; it does not provide the global shadowing, nondegeneracy, or transition-chain estimates needed for diffusion.
\subsection{Short English statement}
The merged source asks for diffusion along every sufficiently small amplitude of one typical perturbation, for a trigonometric-polynomial version, and for a comparison of variational and separatrix-map methods. The aid identifies genericity quantifier gaps, gives a conditional finite-parameter reduction, and records a local generating-function correspondence. None is a supplied proof of the full conjecture.
\subsection{Sources}
\begin{itemize}
\item Local UnsolvedMath registry, permanent ID 20002481, catalogue entry; canonical \texttt{aim-pdes-notes.json}, notes[92], extracted source item 3.
\item AIM workshop list \emph{New connections between dynamical systems and PDE's}, notes by S. Bolotin, \url{https://aimath.org/WWN/dynpde/dynpde.pdf}. The supplied verification identifies the merged entries as Problems 23 and 24, with the leading 2 lost from 23. This draft does not independently verify that account.
\item Source attributions: Mather, Xia; de la Llave, Delshams, Seara; Treschev; Nekhoroshev theory. No paper titles, dates, or theorem references are supplied.
\item Machine-generated research attempt 1, labeled partial progress, proved, and low novelty confidence; not peer reviewed. Its summary reports modern smooth strong diffusion in two and a half degrees of freedom, weaker order-one drift in arbitrary time-periodic dimension, and autonomous cusp-residual same-energy connections, without identifying citations. Its statement that the full fixed-direction and fixed-order versions were not found solved is a report of that attempt, not a verified current literature assessment.
\end{itemize}
\end{document}
